\section{Regla de L'Hôpital}

\begin{teorema}
Sean $f$ y $g$ funciones con valores reales, derivables en un intervalo abierto $I=(a,b)$, tales que $g'(x)\neq 0$ para todo $x\in(a,b)$. Supongamos que $\lim_{x\to a+} f(x)=\lim_{x\to a+} g(x)=0$ y que $\lim_{x\to a+}\frac{f'(x)}{g'(x)}$ existe. Entonces el límite $\lim_{x\to a+}\frac{f(x)}{g(x)}$ existe. Más aún,
\begin{equation}
\lim _{x\to a+}\frac{f'( x)}{g'( x)} =\lim _{x\to a+}\frac{f( x)}{g( x)}.
\end{equation}
\end{teorema}

\begin{prueba}
    Extendemos la definición de $f$ y $g$ hasta el punto $a$ al poner $f(a)=g(a)=0$. Extendidas así, estas funciones son continuas en el intervalo $[a,b)$: la continuidad en todo $x\in(a,b)$ se sigue de la derivabilidad, y la continuidad por la derecha en $a$ se mantiene ya que $\lim_{x\to a+} f(x)=\lim_{x\to a+} g(x)=0$.

    Sea $L=\lim_{x\to a+}\frac{f'(x)}{g'(x)}$. Dado $\varepsilon>0$, existe $\delta>0$, con $\delta < b-a$ (para que $(a,a+\delta)\subset(a,b)$), tal que
    \[
    \left| \frac{f'( t)}{g'( t)} -L\right| < \varepsilon
    \]
    para todo $t\in ( a,a+\delta )$.

    Sea $x\in(a,a+\delta)$. Por el teorema del valor medio generalizado aplicado a $f$ y $g$ en $[a,x]$, existe $\xi \in ( a,x)$ tal que
    \begin{equation}
    \frac{f( x)}{g( x)} =\frac{f( x) -f( a)}{g( x) -g( a)} =\frac{f'( \xi )}{g'( \xi )}.
    \end{equation}
    En efecto: la primera igualdad se cumple porque, tras la extensión, $f(a)=g(a)=0$ (nótese que $g(x)\neq 0$; si $g(x)=0$, el teorema de Rolle aplicado a $g$ en $[a,x]$ produciría un punto de $(a,x)$ con derivada nula, contradiciendo que $g'\neq 0$ en $(a,b)$); y la segunda igualdad es la conclusión del teorema del valor medio generalizado. Por lo tanto, como $\xi\in(a,x)\subset(a,a+\delta)$,
    \[
    \left| \frac{f( x)}{g( x)} -L\right| =\left| \frac{f'( \xi )}{g'( \xi )} -L\right| < \varepsilon.
    \]
    Como $\varepsilon>0$ es arbitrario, entonces $\lim_{x\to a+}\frac{f( x)}{g( x)} =L$.
\end{prueba}

\begin{remark}
     El resultado anterior también es verdadero:
     \begin{itemize}
         \item Si $\lim_{x\to a+} f( x)$ y $\lim_{x\to a+} g( x)$ son ambos $\pm \infty$ en lugar de $0$.
         \item Si los límites se toman en $+\infty$ en vez de $a^+$. En este caso $I = ( b,+\infty )$.
         \item Si los límites se toman en $-\infty$ en vez de $a^+$. En este caso $I = ( -\infty,a )$.
     \end{itemize}
\end{remark}

\begin{prueba}
    \textbf{Caso $\lim_{x\to a+} f( x)=\pm\infty$ y $\lim_{x\to a+} g( x)=+ \infty$.}
    Por hipótesis $\lim_{x\to a+}\frac{f'( x)}{g'( x)} =L$; es decir, para todo $\varepsilon > 0$ existe $\delta_{1} > 0$ tal que si $a< x< a+ \delta_{1}$ entonces
    \[
    L-\frac{\varepsilon }{2} < \frac{f'( x)}{g'( x)} < L+\frac{\varepsilon }{2}.
    \]
    Fijemos un punto $d$ con $a< d< a+\delta_{1}$. Como $f$ y $g$ son derivables en $(a,b)$, son continuas en $[x,d]$ y derivables en $(x,d)$ para cualesquiera $a< x< d$. Por el teorema del valor medio generalizado de Cauchy, existe $\xi \in ( x,d)$ tal que
    \[
    \frac{f'( \xi )}{g'( \xi )} =\frac{f( d) -f( x)}{g( d) -g( x)},
    \]
    y, como $\xi \in (x,d) \subset ( a,a+ \delta _{1})$, se tiene
    \[
    L-\frac{\varepsilon }{2} < \frac{f( d) -f( x)}{g( d) -g( x)} < L+\frac{\varepsilon }{2}.
    \]
    Como $\lim_{x\to a+} g( x) =+ \infty$, existe $\delta_{2} > 0$ tal que si $a< x< a+ \delta _{2}$ entonces $g( x) > g( d)$ y $g( x) > 0$; en particular, el factor $\frac{g( x) -g( d)}{g( x)}$ es positivo. Multiplicando por $\frac{g( x) -g( d)}{g( x)}$ la desigualdad anterior (y observando que $\frac{f( d) -f( x)}{g( d) -g( x)}=\frac{f( x) -f( d)}{g( x) -g( d)}$) obtenemos:
    \begin{equation*}
    \left( L-\frac{\varepsilon }{2}\right)\left(\frac{g( x) -g( d)}{g( x)}\right) < \frac{f( x) -f( d)}{g( x)} < \left( L+\frac{\varepsilon }{2}\right)\left(\frac{g( x) -g( d)}{g( x)}\right).
    \end{equation*}
    Adicionando $\frac{f( d)}{g( x)}$ a los tres miembros, y usando que $\frac{f( x) -f( d)}{g( x)}+\frac{f( d)}{g( x)}=\frac{f( x)}{g( x)}$:
    \begin{equation*}
    L-\frac{\varepsilon }{2} +\frac{f( d) -Lg( d) +\frac{\varepsilon }{2} g( d)}{g( x)} < \frac{f( x)}{g( x)} < L+\frac{\varepsilon }{2} +\frac{f( d) -Lg( d) -\frac{\varepsilon }{2} g( d)}{g( x)}.
    \end{equation*}
    Sea
    \[
    C_1=f( d) -Lg( d) +\frac{\varepsilon }{2} g( d),
    \qquad
    C_2=f( d) -Lg( d) -\frac{\varepsilon }{2} g( d),
    \qquad
    M=\max\{| C_1| ,| C_2| \},
    \]
    que es una constante que no depende de $x$. Entonces, para $a< x< a+\delta _{2}$:
    \begin{equation*}
    \left| \frac{f( x)}{g( x)} -L\right| < \frac{\varepsilon }{2} +\frac{M}{| g( x)| }.
    \end{equation*}
    Como $g( x) \to +\infty$ cuando $x\to a+$, existe $\delta _{3} > 0$ tal que si $a< x< a+ \delta _{3}$ entonces $\frac{M}{| g( x)| } < \frac{\varepsilon }{2}$. Sea $\delta _{m} =\min\{\delta _{2} ,\delta _{3} ,\, d-a\}$. Si $a< x< a+\delta _{m}$ entonces, en particular, $x< d$, de modo que el teorema del valor medio de Cauchy se aplica en $[x,d]$ y todas las condiciones anteriores se cumplen; por lo tanto
    \[
    \left| \frac{f( x)}{g( x)} -L\right| < \frac{\varepsilon }{2} +\frac{\varepsilon }{2} =\varepsilon.
    \]
    Es decir, $\lim _{x\to a+}\frac{f( x)}{g( x)} =L$.
\end{prueba}
